\chapter{PROJECT X}

I want to note that I have not before seen the interpretation of
a typed calculus, taking both the type and typing context as parameters
to give the denotation of a term. This definition of terms with they're typing
context, actually massively simplifies the lean development
previous iterations, involved a type inference algorithm, which could fail
and the denotation function was partial.

Possible goals of this project:
\begin{itemize}
\item Formalisation of Lilac with tactical support
\item Lean interpreter for the language which is sound w.r.t to the
denotational semantics
\item The first case of Compositional Symbolic execution for probabilistic separation logic
  \subitem CSE for the Bluebell separation logic with possible immediate application in industry
  for verifying correctness of Zero Knowledge Proof systems
  This could perhaps be a tactic, which actually invokes the symbolic execution
  to elaborate into a sequence of proof rules?

  This is possibly not so hard
  \subitem CSE for Lilac or BaSL. this is perhaps out of scope of project. But it
  connects to the first part of the project formalising Lilac/BaSL which is novel
  for being the first mechanisation of a separation logic for continuous probabilistic
  programs.
\end{itemize}

terms in
context